\chapter{Classical Mechanics and Dynamics}

\section{Kinematics and Rigid Body Motion}
At the macroscopic scale, Layer 1 interactions are dominated by classical mechanics. Here we define the reference frames, coordinate systems, and the fundamental equations of motion (Newtonian and Lagrangian formulations) that will drive the physics engine.

\section{Collision Detection and Response}
Hard representations require strict boundaries. We must define how objects intersect, transfer momentum, and conserve energy. 

\section{Constraints and Solvers}
Physical constraints (joints, hinges, rigid structural bonds) must be modeled mathematically. This section explores the formulation of Jacobians and the use of iterative solvers (like Gauss-Seidel) to maintain the structural integrity of complex assemblies within the simulation.

\section{Macroscopic Transport and Vector Mechanics}
Beyond microscopic collisions, Layer 1 must model macroscopic kinematics to support logistical simulation (e.g., Scenario Beta: Transit). The simulation engine evaluates the movement of entities (vehicles, humans) as temporal-spatial vectors across the environment grid.

\subsection{Trajectory Overlap and Detour Efficiency}
When attempting to optimize logistics, the simulation evaluates the efficiency of combining entity trajectories. Let $\vec{V}_1$ represent the primary trajectory vector of a host vehicle, and $\vec{V}_2$ represent the requested trajectory of a secondary payload. The orchestrator computes the detour vector $\vec{V}_d$, which is the path required to merge the two. 

The thermodynamic efficiency of the detour is evaluated using standard vector mechanics:
\begin{equation}
\eta_{detour} = \frac{||\vec{V}_2||}{||\vec{V}_1 + \vec{V}_d|| - ||\vec{V}_1||}
\end{equation}
If the additional distance incurred by the detour ($||\vec{V}_1 + \vec{V}_d|| - ||\vec{V}_1||$) significantly outweighs the original payload vector $||\vec{V}_2||$, the simulation rejects the state change as thermodynamically invalid. This rigid vector math prevents the engine from approving highly inefficient logistical routing, forcing macroscopic entity flows to behave like fluid currents seeking the path of least resistance.

\section{Tribology, Wear Rates, and Material Fatigue}
The physical degradation of assets in Layer 1 (e.g., Scenario Nu: Textiles and Soft Infrastructure) is modeled using tribology (the study of friction and wear) and material fatigue.

\subsection{Linear Damage Accumulation}
When an entity engages in physical labor, the kinetic interactions induce micro-structural damage to their equipped tools or soft infrastructure (clothing). The degradation of structural integrity $D$ over time is modeled using the Palmgren-Miner linear damage hypothesis:
\begin{equation}
D = \sum_{i} \frac{n_i}{N_i}
\end{equation}
where $n_i$ is the number of stress cycles applied at a specific stress amplitude, and $N_i$ is the number of cycles to failure at that amplitude (derived from the material's S-N curve). For example, a \texttt{LINEN\_SHIRT} possesses a much steeper S-N curve than a \texttt{CANVAS\_JACKET}.

When $D$ reaches a critical threshold (e.g., $D \ge 0.8$), the material experiences macroscopic failure (a tear). The Oasis engine's MEND mechanic represents a localized injection of negentropy. By applying highly targeted labor and a fraction of the original raw material (thread, patches), the tailor effectively resets $D \to 0$, circumventing the massive thermodynamic cost of fabricating a replacement from virgin polymers or natural fibers \cite{dowling2012mechanical}.



\section{Structural Failure under Wind Load}
Environmental shocks (Scenario Upsilon) introduce massive kinetic forces into the simulation. The engine models the destructive potential of severe weather systems (e.g., high winds acting on biological assets like trees or architectural infrastructure).

\subsection{Aerodynamic Drag and Bending Stress}
The force exerted by wind on a vertical structural voxel (e.g., a \texttt{TREE\_TRUNK}) is modeled using the drag equation:
\begin{equation}
F_d = \frac{1}{2} \rho v^2 C_d A
\end{equation}
where $\rho$ is the density of air, $v$ is wind velocity, $C_d$ is the drag coefficient of the canopy, and $A$ is the projected frontal area. This force creates a bending moment $M = F_d \times h$ at the base of the structure (where $h$ is the center of pressure).

If the resulting bending stress $\sigma = M / Z$ (where $Z$ is the section modulus of the trunk) exceeds the material's modulus of rupture, the physics engine immediately executes a structural failure state transition. The \texttt{TREE\_TRUNK} voxel is converted into a dynamic \texttt{FALLEN\_LOG} entity, which is then subjected to gravity and collision physics, potentially blocking critical pathfinding routes (\texttt{ROAD} voxels) and necessitating immediate kinetic labor from the community to resolve \cite{gordon2003structures}.

\begin{thebibliography}{99}
\bibitem{dowling2012mechanical}
Dowling, N. E. (2012). \textit{Mechanical behavior of materials}. Pearson.

\bibitem{gordon2003structures}
Gordon, J. E. (2003). \textit{Structures: or why things don't fall down}. Da Capo Press.

\end{thebibliography}
